\documentclass{article}
\usepackage{../../../math_template}

\author{Jonathan Kasongo}
\title{HW 2}

\begin{document}
\maketitle

Problems from: \url{https://www.math.cmu.edu/~ellenp/HWK/ma355hw2.pdf}

\textit{Due: 5/6/26}

\begin{problem}{Problem 1}
Let $A, B, C$ be sets and let $f : A \rightarrow B$, $g : B \rightarrow C$
be functions. Prove: If $f$ is onto $B$ and $g$ is onto $C$, then
$g \circ f : A \rightarrow C$ is onto $C$.
\end{problem}

\begin{solution}{Solution}
Since $f$ is surjective, we have
$\forall b \in B, \ \exists a \in A \text{ such that } f(a) = b$. \\

Since $g$ is surjective, we have
$\forall c \in C, \ \exists b \in B \text{ such that } f(b) = c$. \\

Let $c_0 \in C$ be arbitrary
then $\exists b_0 \in b$ such that $g(b_0) = c_0$, by the second statement.
Similarly $\exists a_0 \in a$ such that $f(a_0) = b_0$, by the first
statement. So $g(f(a_0)) = c_0$. So we have shown that for any $c_0 \in C$,
$\exists a_0 \in A$ such that $g(f(a_0)) = c_0$, and so $g\circ f$ is
surjective. $\blacksquare$
\end{solution}
\vspace{0.4cm}

\begin{problem}{Problem 2}
Show that the relation $\sim$ (two sets are equivalent) is an equivalence
relation.
\end{problem}

\begin{solution}{Solution}
We need three properties. Reflexivity, symmetry and transitivity. Let
$A,B,C$ be sets.
\\

Trivially $A \sim A$. So $\sim$ is reflexive. \\

Clearly if $A$ and $B$ are the same set then $A \sim B$ and $B \sim A$.
So $\sim$ is symmetric. \\

If $A$ and $B$ are the same set and $B$ and $C$ are the same set then
$A$ and $C$ are the same set, so it means that
$A \sim B \text{ and } B \sim C \Rightarrow A \sim C$. So $\sim$ is
transistive. $\blacksquare$
\end{solution}
\vspace{0.4cm}

\begin{problem}{Problem 3}
Give an example of a countable collection of finite sets whose union is
not finite.
\end{problem}

\begin{solution}{Solution}
For example the collection of sets
$$
\{ 1 \}, \{ 2 \}, \{ 3 \}, \ldots
$$
is countable and it's union is $\mathbb{N}$ which is infinite.
$\blacksquare$
\end{solution}
\newpage

\begin{problem}{Problem 4}
Are the following sets finite, countable or uncountable? Explain or
prove your answer in each case
\begin{enumerate}
\item $\{ (x,y) \in \mathbb{N} \times \mathbb{R} : xy = 1 \}$

\item $(\frac14, \frac34)$
\end{enumerate}
\end{problem}

\begin{solution}{Solution}
1. Note that $y$ is uniquely determined by $x$ so the set is equivalent to
$$
\left\{ \left(x, \frac{1}{x}\right) : x \in \mathbb{N} \right\}
$$
which is clearly countable. \\

2. This interval is uncountable. Let us show it by contradiction.
Suppose that $f$ is a bijection mapping
$\mathbb{N} \rightarrow (\frac14, \frac34)$. We may assume that
$f$ is strictly increasing, since we can re-order the elements in the
image of $f$ as we wish. Define, for some fixed $k \in \mathbb{N}$
$$
d = \frac12 \left( f(k+1) + f(k) \right).
$$
We have $f(k) < d < f(k+1)$ so $d \in (\frac14, \frac34)$ but $d$ isn't in
the image of $f$, a contradiction. $\blacksquare$
\end{solution}
\vspace{0.4cm}

\begin{problem}{Problem 5}
Is the set of all irrational numbers countable? Prove your answer.
\end{problem}

\begin{solution}{Solution}
The answer is no, the set of irrational numbers is uncountable. Define
$I$ to be the set of irrational numbers. Suppose
FTSOC that $I$ is countable. Then any non-empty subset of $I$ is countable.
\\

Let $a < b$ be real. Take $S(a,b) := \{ x : x \in I, \ a \leq x \leq b \}$.
Let $R(a,b) := \{ x : x \in \mathbb{Q}, \ a \leq x \leq b \}$. It is
well-known that $\mathbb{Q}$ is countable and since
$R(a,b) \subset \mathbb{Q}$, $R(a,b)$ is countable. Observe that
$$
S(a,b) \cup R(a,b) = [a,b]
$$
is countable since the union of countable sets is countable. However any
non-empty interval of real numbers is uncountable, contradiction.
$\blacksquare$
\end{solution}
\newpage

\begin{problem}{Problem 6}
Prove that $\mathbb{N} \times \mathbb{N}$ is countable.

\textit{Hint: Consider } $f(m,n) = 2^{m-1}(2n-1)$.
\end{problem}

\begin{solution}{Solution}
The idea is to show that
$f : \mathbb{N} \times \mathbb{N} \rightarrow \mathbb{N}$ is a bijection.
\\

Note that for any $n \in \mathbb{N}$ we can write
$$
n = 2^{\nu_2(n)} \frac{n}{2^{\nu_2(n)}}
$$
where $\nu_p$ is the $p$-adic valuation. We mean to say that any natural
number can be expressed as a pure power of 2 multiplied by some odd number.
Since $m \geq 1$, $2^{m-1}$ can be any power of 2 in $\mathbb{N}$,
and $2n-1$ can be any odd number in $\mathbb{N}$. Hence $f$ is surjective.
\\

Next we show that $f$ is injective. Suppose that $f(a,b) = f(m,n)$.
Clearly $a = m$ because otherwise $\nu_2(f(a,b)) \neq \nu_2(f(m,n))$.
Thus we have
$$
2n-1 = 2b-1 \iff n=b
$$
and so $f$ is injective. Hence $\mathbb{N} \times \mathbb{N}$ is in
bijection with $\mathbb{N}$ and so it is countable. $\blacksquare$
\end{solution}

\end{document}
